\documentclass[UTF8,11pt,fontset=none]{ctexart}
\usepackage[a4paper,margin=23mm,headheight=15pt,footskip=11mm]{geometry}
\usepackage{fontspec}
\IfFontExistsTF{TeX Gyre Termes}{\setmainfont{TeX Gyre Termes}}{\setmainfont{lmroman10-regular.otf}[BoldFont=lmroman10-bold.otf,ItalicFont=lmroman10-italic.otf,BoldItalicFont=lmroman10-bolditalic.otf]}
\IfFontExistsTF{TeX Gyre Heros}{\setsansfont{TeX Gyre Heros}}{\setsansfont{lmsans10-regular.otf}[BoldFont=lmsans10-bold.otf,ItalicFont=lmsans10-oblique.otf]}
\IfFontExistsTF{DejaVu Sans Mono}{\setmonofont{DejaVu Sans Mono}[Scale=MatchLowercase]}{\setmonofont{lmmono10-regular.otf}}
\IfFontExistsTF{Noto Serif CJK SC}{\setCJKmainfont{Noto Serif CJK SC}[AutoFakeSlant=0.15]}{\setCJKmainfont{FandolSong-Regular.otf}[BoldFont=FandolSong-Bold.otf,ItalicFont=FandolKai-Regular.otf]}
\IfFontExistsTF{Noto Sans CJK SC}{\setCJKsansfont{Noto Sans CJK SC}}{\setCJKsansfont{FandolHei-Regular.otf}}
\IfFontExistsTF{Noto Sans Mono CJK SC}{\setCJKmonofont{Noto Sans Mono CJK SC}}{\setCJKmonofont{FandolFang-Regular.otf}}
\usepackage{amsmath,amssymb,amsthm,mathtools,mathrsfs,esint,bm}
\usepackage{booktabs,longtable,array,calc}
\usepackage{enumitem}
\setlist{nosep,leftmargin=1.7em,topsep=3pt}
\usepackage{fvextra}
\DefineVerbatimEnvironment{Highlighting}{Verbatim}{breaklines,breakanywhere,fontsize=\small}
\usepackage{xurl}
\usepackage[unicode,hidelinks,bookmarksnumbered]{hyperref}
\usepackage{bookmark}
\usepackage{fancyhdr}
\pagestyle{fancy}\fancyhf{}
\fancyhead[L]{\small 第一科学主线：近期研究全稿}
\fancyhead[R]{\small 2026-10-07}\fancyfoot[C]{\thepage}
\renewcommand{\headrulewidth}{0.3pt}
\providecommand{\tightlist}{\setlength{\itemsep}{1pt}\setlength{\parskip}{0pt}}
\providecommand{\passthrough}[1]{#1}
\newcommand{\Ric}{\operatorname{Ric}}
\newcommand{\Vol}{\operatorname{Vol}}
\newcommand{\Rm}{\operatorname{Rm}}
\newcommand{\supp}{\operatorname{supp}}
\newcommand{\Var}{\operatorname{Var}}
\newcommand{\tr}{\operatorname{tr}}
\newcommand{\avint}{\fint}
\newcommand{\Jcal}{\mathcal J}
\newcommand{\Ecal}{\mathcal E}
\newcommand{\dd}{\,\mathrm d}
\newcommand{\Rea}{\operatorname{Re}}
\setlength{\parindent}{2em}
\setlength{\parskip}{3pt plus 1pt}
\setlength{\emergencystretch}{3em}
\linespread{1.10}\allowdisplaybreaks[2]
\setlength{\abovedisplayskip}{8pt plus 2pt minus 2pt}
\setlength{\belowdisplayskip}{8pt plus 2pt minus 2pt}
\renewcommand{\arraystretch}{1.12}
\setcounter{tocdepth}{1}\setcounter{secnumdepth}{2}
\pdfstringdefDisableCommands{\def\fint{int}\def\mathcal#1{#1}\def\mathfrak#1{#1}\def\mathscr#1{#1}\def\operatorname#1{#1}\def\varepsilon{epsilon}\def\chi{chi}\def\beta{beta}\def\rho{rho}\def\tau{tau}\def\Phi{Phi}\def\sigma{sigma}\def\kappa{kappa}\def\nabla{D}\def\mathbb#1{#1}\def\({}\def\){}\def\textsuperscript#1{#1}\def\textsubscript#1{#1}}
\title{第五幂线丛与任意固定有限密度端点}
\author{第一科学主线·近期研究分卷}
\date{2026年10月7日}
\begin{document}
\maketitle
\noindent\textbf{原稿编号：}R03。

本章收录本会话的第五幂线丛证明。直接输入包括强流、有界同丛 resolvent，以及所引第一延拓；结论是在任意固定有限 Euler 密度下的端点估计。其他网页随后给出的秩一陈数证明是同一结论的另一条路线，作为直接输入单独登记。\cite{R03,Strong,Prolongation,FD}

\section{有限密度结论}\label{ch3:ux6709ux9650ux5bc6ux5ea6ux7ed3ux8bba}

在下文直接前置成立时，不存在 \[
s_j=-4,\quad\beta_j\to\infty,\quad
\kappa_j=(P_j+D_{\beta_j})/V_j\to0,
\qquad \sup_j\fint|W_j^-|^2<\infty
\tag{3.1}\label{eq:r3:1}
\] 的原允许闭流形序列。因此，对每个固定有限 \(B\ge0\)，存在 \(c(B)>0\)，使 \[
\frac{\chi(M)}A\le B\quad\Longrightarrow\quad
\kappa_\beta(g)\ge c(B)\frac{(\beta-2)^2}{\beta^2},\qquad\beta\ge2.
\tag{3.2}\label{eq:r3:2}
\] 这里 \(p\ge A>0\)、\(a^2\ge T_+\) 及实际共形供给保持原条件。

\section{固定范数和直接前置}\label{ch3:ux56faux5b9aux8303ux6570ux548cux76f4ux63a5ux524dux7f6e}

记 \(L=\det S^+\)、\(E=S^-\otimes L\)，故 \(\det E=L^3\)。\(A\) 是 \(L\) 的实际 unitary 联络。为避免拓扑预算同名，局部写 \(A=i\mathfrak a\)，\(\mathfrak a\) 为实一形式；其曲率为 \(i\,d\mathfrak a\)。线丛 \(L^k\) 的导数记 \(D_k\)，局部为 \(d+ki\mathfrak a\)。所有 rough Laplacian 取 \(\nabla^*\nabla\ge0\)。Hermitian 内积对第一变量共轭线性，最后取实部。

二形式采用算子／外代数规范，\(|e^{12}\pm e^{34}|^2=2\)。实际负手性二次映射满足 \[
 c(q_-(z))z=\frac{|z|^2}{2}z,\qquad |q_-(z)|^2=\frac{|z|^4}{8}.
\]

假设存在 \eqref{eq:r3:1}。原共形归约及同解估计提供 \[
 D_A\Phi=0,\quad F_A^+=2i\sigma(\Phi),\quad
 \rho=|\Phi|^2\le2,\qquad \rho_* =4/3,
\] \[
 \tau=t_\Phi\in\Gamma(E),\quad m_*=128/81,
 \quad |\tau|^2\to m_*\text{ in }L^2_{\rm av},
 \quad D^-_{3A}\tau\to0\text{ in }L^2_{\rm av}.
\] 正旋量具有 \(\nabla\Phi\) 的归一化 \(L^4\) 一致可积性及 \(\nabla^2\Phi\) 的归一化 \(L^2\) 界，且 \[
 \rho\to\rho_*,\quad \nabla\zeta\to0,\quad
 \zeta=\langle\Phi,\nabla_A\Phi\rangle\to0,
 \quad N_\tau=(\nabla\tau)^{1,0}_{J_+}\to0.
 \tag{3.3}\label{eq:r3:3}
\] 这里所用小量均取其已证的归一化 \(L^2\) 范数；\(\zeta\) 另由一致可积性趋零于 \(L^4\)。

续研 B 的有界截断与同丛 resolvent 给出 \(v\in\Gamma(E)\)，满足 \[
 |v|+|\nabla^*\nabla v|\le C,
 \quad v-\tau\to0\text{ in }H^1_{\rm av}\cap L^4_{\rm av},
\] \[
 D^-_{3A}v\to0,\quad
 (\nabla^*\nabla-2)v\to0\text{ in }L^2_{\rm av},\quad
 m:=|v|^2\to m_*\text{ in }L^2_{\rm av}.
 \tag{3.4}\label{eq:r3:4}
\] 由 \eqref{eq:r3:1} 的有限平均 \(W^-\) 预算和已证 Hessian 交换， \[
 \|\nabla v\|_{4,\rm av}+\|\nabla^2v\|_{2,\rm av}\le C_K,
 \qquad \fint|\nabla v|^2\to2m_*.
 \tag{3.5}\label{eq:r3:5}
\] 其中 \(K=\sup_j\fint |W^-|^2<\infty\)。另有 \[
 dm\to0,\qquad j_v:=\operatorname{Im}\langle v,\nabla v\rangle\to0
 \quad\text{in }L^2_{\rm av}.
 \tag{3.6}\label{eq:r3:6}
\]

\textbf{第一延拓是本文必需的直接输入。} 在 \(\rho,m\) 远离零的集合上，取后述适配局部标架。记 \(Z_1=e_1-ie_2\)、\(Z_2=e_3-ie_4\)。续研 B 的 \path{O10_Resolvent_Prolongation.md} §§2--4 证明 \[
 \nabla_{Z_1}v,\quad \nabla_{Z_2}v,\quad
 \nabla_{\bar Z_1}v\to0\quad\text{in }L^2_{\rm av};
 \tag{3.7}\label{eq:r3:7}
\] Dirac 迹再使 \((\nabla_{\bar Z_2}v)_1\to0\)。只有 \((\nabla_{\bar Z_2}v)_2\) 未被消去。第一延拓仅用有限 \(K\) 的导数预算，并不使用 \(K\le50/3\) 或临界谱等号。

低 \(\rho\)、低 \(m\) 集合的相对体积趋零。由 \eqref{eq:r3:5}，这些集合上的 \(\nabla v\) 二次能量为 \(o(V)\)。以后所有适配标架计算先作固定幅度截断，再延拓为全局张量；不假设存在全局标架。

同一实际曲率拟合还给出 \[
 f:=-iF_A^-=-\frac{27}{32}Q_v+o_{L^2_{\rm av}}(1),
 \qquad Q_v=-iq_-(v).
 \tag{3.8}\label{eq:r3:8}
\] 这是原有 \(f-K_{\rm fit}\to0\)、\(K_{\rm fit}=-(9/8)\rho^3Q_\tau/S\)，以及 \(v-\tau\to0\) 于 \(L^4\) 的直接结果；极限系数为 \((9/8)\rho_*^3/(2m_*)=27/32\)。不求这个拟合模型的导数。

\section{\texorpdfstring{剩余导数的真实载体是线丛 \(L^5\)}{剩余导数的真实载体是线丛 L\^{}5}}\label{ch3:ux5269ux4f59ux5bfcux6570ux7684ux771fux5b9eux8f7dux4f53ux662fux7ebfux4e1b-l5}

\subsection{1 适配标架和误差}\label{ch3:ux9002ux914dux6807ux67b6ux548cux8befux5dee}

在 \(\rho,m>0\) 上，选适配 Spin\(^c\) 和正交标架，使 \[
 \Phi=\sqrt\rho\,(1,0),\qquad v=\sqrt m\,(1,0).
\] 用 \[
 \theta_1=e^1-ie^2,\qquad \theta_2=e^3-ie^4
\] 记复一形式。正、负扭曲旋量联络可写为 \[
 \Omega_+=\begin{pmatrix}ip&-\bar\eta\\\eta&i(\mathfrak a-p)\end{pmatrix},
 \qquad
 \Omega_E=\begin{pmatrix}iq&-\bar b\\b&i(3\mathfrak a-q)\end{pmatrix},
 \tag{3.9}\label{eq:r3:9}
\] 其中 \(p=\operatorname{Im}\zeta/\rho\)、\(q=j_v/m\) 是实相位流，不是单极子投影平方。

实际三次式、\(v-\tau\to0\) 及小流给 \[
 \eta=\kappa\theta_1+\epsilon_\eta,
 \qquad \kappa=\frac{\sqrt m}{2\rho^{3/2}}\in\mathbb R,
 \quad \epsilon_\eta\to0\text{ in }L^2_{\rm av}.
 \tag{3.10}\label{eq:r3:10}
\] 第一延拓给 \[
 b=c\theta_2+\epsilon_b,\qquad
 c=\frac{b_3+ib_4}{2},\qquad \epsilon_b\to0\text{ in }L^2_{\rm av}.
 \tag{3.11}\label{eq:r3:11}
\] 在任一固定下幅度截断上，所有 \(p,q,\eta,b,c,\epsilon_\eta,\epsilon_b\) 具有归一化 \(L^4\) 界，且 \(p,q\to0\) 于 \(L^2\)。这些界来自 \(\nabla\Phi,\nabla v\) 的界，不涉及裸的局部势 \(\mathfrak a\) 的界。

\subsection{2 不依赖标架的定义}\label{ch3:ux4e0dux4f9dux8d56ux6807ux67b6ux7684ux5b9aux4e49}

对 \(v\in E\)，定义其三次表达式 \[
 t(v)=\sum_{i=1}^4 c^-(e^i)v\otimes(v\wedge\nabla_i v)
       \in\Gamma(S^+\otimes L^4).
\] 因此 \[
 \Phi\wedge t(v)\in\Gamma(L^5).
\] 在上述标架和稿件 Clifford 规范下， \[
 t(v)=m^{3/2}(-b_1+ib_2,-b_3-ib_4),
\] 所以 \[
 c_0=-\frac{\Phi\wedge t(v)}{2\sqrt\rho\,m^{3/2}}
 \tag{3.12}\label{eq:r3:12}
\] 的局部系数正好是 \eqref{eq:r3:11} 中的 \(c\)。

选固定光滑截断 \(\chi_\rho,\chi_m\)，在零附近为零，在 \(\rho_*\)、\(m_*\) 的固定邻域内为一；其导数支撑与这两个极限幅度保持正距离。定义 \[
 c_j=\chi_\rho(\rho_j)\chi_m(m_j)c_{0,j}\in\Gamma(L_j^5),
 \tag{3.13}\label{eq:r3:13}
\] 并在截断为零的地方光滑延拓。以下全局 \(c\) 始终指 \eqref{eq:r3:13}。

若局部行列式标架变为 \(e^{it}\) 倍，\(c_0\) 的系数变为 \(e^{-5it}c_0\)，\(\theta_2\) 变为 \(e^{2it}\theta_2\)，\(b\) 变为 \(e^{-3it}b\)。因此 \eqref{eq:r3:11} 与线丛连接 \(D_5=d+5i\mathfrak a\) 相容。

\subsection{3 导数预算和非消失质量}\label{ch3:ux5bfcux6570ux9884ux7b97ux548cux975eux6d88ux5931ux8d28ux91cf}

对截断区间，直接对 \eqref{eq:r3:12} 求协变导数，有 \[
 |c|\le C|\nabla v|,
\] \[
 |D_5c|\le C\left(|\nabla^2v|+|\nabla v|^2
                         +|\nabla\Phi|\,|\nabla v|\right).
 \tag{3.14}\label{eq:r3:14}
\] 截断导数属于右侧同一类型。因此 \[
 \|c_j\|_{4,\rm av}+\|D_5c_j\|_{2,\rm av}\le C_K.
 \tag{3.15}\label{eq:r3:15}
\]

更重要的是 \[
 \boxed{\fint|c_j|^2\longrightarrow1.}
 \tag{3.16}\label{eq:r3:16}
\] 证明：在好幅度区， \[
 |\nabla v|^2=\frac{|dm|^2}{4m}+m|q|^2+m|b|^2.
\] 前两项为 \(o(V)\)；\eqref{eq:r3:11} 给 \(|b|^2=2|c_0|^2+o_{L^1_{\rm av}}(1)\)。由于 \(m-m_*\to0\) 于 \(L^2\)、\(|b|^2\) 在 \(L^2\) 中有界，可以将乘子 \(m\) 替换为 \(m_*\)。坏幅度区的梯度二次能量由 \eqref{eq:r3:5} 为 \(o(V)\)。结合 \(\fint|\nabla v|^2\to2m_*\) 即得 \eqref{eq:r3:16}。这个步骤不要求 \(|\nabla v|^4\) 一致可积，只要求其有界。

\section{连接剩余分量：带两条相位流的精确 Cartan 计算}\label{ch3:ux8fdeux63a5ux5269ux4f59ux5206ux91cfux5e26ux4e24ux6761ux76f8ux4f4dux6d41ux7684ux7cbeux786e-cartan-ux8ba1ux7b97}

采用原稿 \(S^+\to S^-\) Clifford 矩阵 \[
 C_1=I,\quad C_2=\operatorname{diag}(i,-i),\quad
 C_3=\begin{pmatrix}0&1\\-1&0\end{pmatrix},\quad
 C_4=\begin{pmatrix}0&-i\\-i&0\end{pmatrix};
\] 反向矩阵是 \(-C_i^*\)。LC 的两手性 traceless 连接分别从 \eqref{eq:r3:9} 去掉 \(i\mathfrak a I/2\)、\(3i\mathfrak a I/2\)。由 \[
 \Omega_{-,\mathrm{LC}}C_j-C_j\Omega_{+,\mathrm{LC}}
       =\sum_i\Gamma^i{}_j C_i,
 \qquad de^i=-\Gamma^i{}_j\wedge e^j
\] 得到：当 \(\eta=\kappa\theta_1\)、\(b=c\theta_2\) 时， \[
 \boxed{d\theta_2=i(2\mathfrak a-p-q)\wedge\theta_2
                    -\bar c\,\Xi_-,}
 \tag{3.17}\label{eq:r3:17}
\] 其中 \[
 \Xi_-=e^{13}+e^{24}+i(e^{14}-e^{23}),\qquad (\Xi_-)^+=0.
\] 系数 \(\kappa\) 在此式中抵消，不对其求导。对真实的 \eqref{eq:r3:10}--\eqref{eq:r3:11}，同一线性 Cartan 计算给 \[
 d_{-2A}\theta_2=-i(p+q)\wedge\theta_2-\bar c_0\Xi_-+\mathcal R,
 \quad |\mathcal R|\le C(|\epsilon_\eta|+|\epsilon_b|).
 \tag{3.18}\label{eq:r3:18}
\] 这里只是连接的一阶代数关系，不涉及曲率求导。

另一方面，\(E\) 的 off-diagonal 曲率为 \[
 (F^E)_{21}=db+i(3\mathfrak a-2q)\wedge b.
\] 代入 \eqref{eq:r3:11}、\eqref{eq:r3:18}，得到 \[
 \begin{aligned}
 (F^E)_{21}={}&D_5c_0\wedge\theta_2-|c_0|^2\Xi_-
       +(d_{3A}\epsilon_b)\\
 &-i(p+3q)\wedge c_0\theta_2
       -2iq\wedge\epsilon_b+c_0\mathcal R.
 \end{aligned}
 \tag{3.19}\label{eq:r3:19}
\] 取相对于原取向的正自对偶部分。\(|c_0|^2\Xi_-\) \textbf{精确消失}。而 \(((F^E)_{21})^+\) 的中心部分是对角的，非中心部分来自 LC 的 Ricci off-diagonal 块，故 \[
 |((F^E)_{21})^+|\le C|\operatorname{Ric}^{\circ}|.
\] 记 \[
 a_0\mapsto\mathcal L(a_0):=(a_0\wedge\theta_2)^+.
\] 于是 \eqref{eq:r3:19} 给出 \[
 \mathcal L(D_5c_0)=-(d_{3A}\epsilon_b)^++G_j,
 \qquad \|G_j\|_{4/3,\rm av}\to0.
 \tag{3.20}\label{eq:r3:20}
\] 例如 \(pc_0,qc_0\) 由小 \(L^2\) 乘有界 \(L^4\) 控制，\(c_0\mathcal R\) 同样；Ricci 项甚至在 \(L^2\) 中趋零。没有删除 \(W^-\)：它根本不进入这一 off-diagonal 自对偶曲率方程。

\section{把局部 CR 约束变成全局弱测试}\label{ch3:ux628aux5c40ux90e8-cr-ux7ea6ux675fux53d8ux6210ux5168ux5c40ux5f31ux6d4bux8bd5}

令 \(J_-\) 的基本形式为 \(\Omega_-=e^{12}-e^{34}\)，并在一形式上约定 \((J_-\alpha)(X)=-\alpha(J_-X)\)。

对 \(\alpha=\sum\alpha_i e^i\)，\(\mathcal L(\alpha)=0\) 等价于 \[
 \alpha_1+i\alpha_2=0,\qquad\alpha_3-i\alpha_4=0.
\] 它控制的是此 \(J_-\) 的反全纯导数，不是原 \(J_+\) 的同一个方向。

为准确保留权重，取 \(\Psi_1=e^{12}+e^{34}\)、\(\Psi_2=e^{13}-e^{24}\)、\(\Psi_3=e^{14}+e^{23}\)，各自范数平方为 2。\(\mathcal L\) 的坐标为 \[
 \frac12\bigl(-i\alpha_3-\alpha_4,\ \alpha_1+i\alpha_2,
                                      -i\alpha_1+\alpha_2\bigr).
\] 在这个正自对偶三维目标中令 \(\mathcal D=\operatorname{diag}(2,1,1)\)。则精确 Hermitian 符号恒等式为 \[
 \boxed{\mathcal L^*\mathcal D\mathcal L=I-iJ_-.}
 \tag{3.21}\label{eq:r3:21}
\] 这是一个由实际 \(\Phi,v\) 决定的有界 bundle map；在固定下幅度区，其协变一阶导数由 \(C(|\nabla\Phi|+|\nabla v|)\) 控制。

\textbf{弱测试结论。} 对任意光滑 \(w_j\in\Gamma(L_j^5)\)，若 \[
 \|w_j\|_\infty+\|D_5w_j\|_{4,\rm av}
                      +\|D_5^2w_j\|_{2,\rm av}\le C_w,
 \tag{3.22}\label{eq:r3:22}
\] 则 \[
 \Re\fint\langle D_5c_j-i\bar J_jD_5c_j,D_5w_j\rangle\to0.
 \tag{3.23}\label{eq:r3:23}
\] 这里使用全局光滑有界的二形式及其 skew map \[
 \bar\Omega_j=-\frac4{m_*}Q_{v_j},\qquad
 (\bar J_j\alpha)(X)=-\alpha(\bar J_jX).
 \tag{3.24}\label{eq:r3:24}
\] \(\bar J_j\) 不必处处平方等于 \(-I\)，也不必是一个新的 almost-complex 结构。它只在后续积分恒等式中作为 skew tensor 使用。

证明 \eqref{eq:r3:23} 的所有项如下。

\begin{enumerate}
\def\labelenumi{\arabic{enumi}.}
\tightlist
\item
  在固定幅度截断内，以 \(\mathcal D\mathcal L(D_5w)\) 测试 \eqref{eq:r3:20}，并用 \eqref{eq:r3:21}。\(G_j\) 与 \(D_5w\) 的配对由 \(L^{4/3}\)-\(L^4\) 为 \(o(1)\)。
\item
  将 \(d_{3A}\epsilon_b\) 的导数转到测试上。其 adjoint 由 \[
   C\bigl(|D_5^2w|+(|\nabla\Phi|+|\nabla v|+|d\chi|)|D_5w|\bigr)
  \] 控制，因而在 \(L^2_{\rm av}\) 中有界；与 \(\epsilon_b\to0\) 的 \(L^2\) 配对为 \(o(1)\)。这里使用 unitary 的自然 bundle covariant derivatives，局部 \(\mathfrak a\) 项相互匹配，不要求 gauge 势的额外界。
\item
  将 \(c_0\) 换成截断 \(c_j\)，出现 \(c_0d\chi\)。其支撑相对体积趋零；\(c_0,d\chi\) 均 \(L^4\) 有界，所以该项 \(L^2\) 有界并在 \(L^{4/3}\) 中趋零。可以采用两个嵌套的固定下幅度截断，使每个局部计算及测试在零幅度之前光滑为零。
\item
  在好幅度区，\(Q_v=-m\Omega_-/4\)，故 \(\bar J=(m/m_*)J_-\)。\(m\) 有界并趋于 \(m_*\)，所以 \(\bar J-J_-\to0\) 于 \(L^4\)。用 \eqref{eq:r3:15} 及 \eqref{eq:r3:22}，更换 \(J_-\) 的误差为 \(o(1)\)。
\end{enumerate}

这些步骤只用一阶结构误差的小 \(L^2\)、两个 \(L^4\) 因子的乘积，以及测试的二阶 \(L^2\)。没有给 \(\epsilon_b\) 的导数假定强收敛，也没有把 \(H^{-1}\) 残差升级成未经证明的强 \(L^2\) 方程。

\section{同线丛 Bochner 配对给出确定的正势}\label{ch3:ux540cux7ebfux4e1b-bochner-ux914dux5bf9ux7ed9ux51faux786eux5b9aux7684ux6b63ux52bf}

由 Hopf 散度式与 \eqref{eq:r3:4}、\eqref{eq:r3:6}， \[
 \delta\bar\Omega\to0\quad\text{in }L^2_{\rm av}.
 \tag{3.25}\label{eq:r3:25}
\] 线丛 \(L^5\) 的曲率是 \(F^{L^5}=5F_A\)。由 \eqref{eq:r3:8}、\(|Q_v|^2=m^2/8\)， \[
 \begin{aligned}
 \mathfrak q_j&:=-i\langle F^{L^5},\bar\Omega\rangle
       =5\langle f,\bar\Omega\rangle\\
 &\longrightarrow
 5\frac{27}{32}\frac4{m_*}\frac{m_*^2}{8}
   =\frac{135m_*}{64}=\frac{10}{3}
 \quad\text{in }L^2_{\rm av}.
 \end{aligned}
 \tag{3.26}\label{eq:r3:26}
\] 这里正手性 \(F_A^+\) 与 \(\bar\Omega\) 正交。符号由原负手性 Clifford 约定决定，不与第二次三次旋量的 \(-7\) Lichnerowicz 势混同。

对任意全局二形式 \(\Omega\)、同线丛截面 \(c,w\)，直接分部积分及线丛曲率交换给 \[
 \begin{aligned}
 &\Re\int\langle D_5c-iJ_\Omega D_5c,D_5w\rangle\\
 &=\Re\int\left\langle c,\,
    \Delta_{5A}w-i\langle F^{L^5},\Omega\rangle w
                         +i\,\delta\Omega\cdot D_5w\right\rangle.
 \end{aligned}
 \tag{3.27}\label{eq:r3:27}
\] 不需要 \(\Omega\) 闭，也不需要 \(J_\Omega^2=-I\)。在正交标架上，skew 项是 \(i\Omega_{ij}\overline{D_i c}D_jw\)；换导数产生 \(-i\langle F^{L^5},\Omega\rangle\)，其余是 \(i\delta\Omega\cdot Dw\)。

把 \eqref{eq:r3:25}、\eqref{eq:r3:26} 代入 \eqref{eq:r3:27}。势误差与 \(c\)、有界 \(w\) 用 \(L^2\)-\(L^2\) 配对；漂移误差用 \(\|\delta\bar\Omega\|_2\|c\|_4\|D_5w\|_4\)。结合 \eqref{eq:r3:23}，得到 \[
 \boxed{\Re\fint\langle c_j,(\Delta_{5A}+10/3)w_j\rangle\to0}
 \tag{3.28}\label{eq:r3:28}
\] 对所有 \eqref{eq:r3:22} 的测试成立。

这不是 \(c_j\) 的强二阶方程；它是一个足够用来做以下同线丛 resolvent 对偶测试的弱结论。

\section{一个合法的对偶测试消去全部剩余质量}\label{ch3:ux4e00ux4e2aux5408ux6cd5ux7684ux5bf9ux5076ux6d4bux8bd5ux6d88ux53bbux5168ux90e8ux5269ux4f59ux8d28ux91cf}

固定大数 \(M\)，取光滑径向截断 \(T_M\)，在 \(|c|\le M\) 上 \(T_M(c)=c\)，全局 \(|T_M(c)|\le2M\)，且 \[
 \Re\langle c,T_M(c)\rangle\ge0.
\] 在 \(L^5\) 上解 \[
 \boxed{(\Delta_{5A}+10/3)w_j=T_M(c_j).}
 \tag{3.29}\label{eq:r3:29}
\] 正势给可逆性和 \(\|w_j\|_\infty\le3(2M)/10\)，继而 \(\Delta_{5A}w_j\) 有界。

本测试的连接曲率只有 \(5F_A\)，没有 \(W^-\) 块。原实际拟合使它成为有界模型加归一化 \(L^2\) 小误差，Ricci 同样如此。因而已证的有界强迫项 Hessian／有界截面插值给 \[
 \|D_5^2w_j\|_{2,\rm av}+\|D_5w_j\|_{4,\rm av}\le C_M.
\] 这正好验证 \eqref{eq:r3:22}，不是预设它。证明不需要 \(T_M(c_j)\) 的导数，也不需要 \(w_j\) 的 \(L^\infty\) 梯度界。

由 \eqref{eq:r3:28}、\eqref{eq:r3:29}， \[
 \Re\fint\langle c_j,T_M(c_j)\rangle\longrightarrow0.
 \tag{3.30}\label{eq:r3:30}
\] 另一方面，\eqref{eq:r3:15}--\eqref{eq:r3:16} 给 \[
 \liminf_j\Re\fint\langle c_j,T_M(c_j)\rangle
 \ge 1-\limsup_j\fint_{\{|c_j|>M\}}|c_j|^2
 \ge1-\frac{C_K}{M^2}.
 \tag{3.31}\label{eq:r3:31}
\] 先固定足够大的 \(M\)，使右端至少 \(1/2\)，与 \eqref{eq:r3:30} 矛盾。

\textbf{这排除了 \eqref{eq:r3:1}。}关键是测试在 \(L^5\) 上，其连接曲率受原实际行列式预算控制；无需把相反 Weyl 集中的陈类系数优化到锐值。

\section{从序列排除到每一个有限密度范围的统一估计}\label{ch3:ux4eceux5e8fux5217ux6392ux9664ux5230ux6bcfux4e00ux4e2aux6709ux9650ux5bc6ux5ea6ux8303ux56f4ux7684ux7edfux4e00ux4f30ux8ba1}

固定 \(B<\infty\)。若 \eqref{eq:r3:2} 不成立，取一列原输入，使 \[
 \frac{\kappa_{\beta_j}}{(\beta_j-2)^2/\beta_j^2}\to0.
\] \(\beta=2\) 处右侧为零，无需纳入这条反证子列。已有有限参数估计排除任何有界 \(\beta_j\) 子列，因此 \(\beta_j\to\infty\)、\(\kappa_{\beta_j}\to0\)。

已有精确共形比较 \[
 \widehat\kappa\le
 \frac{\kappa}{1-\sqrt\kappa-\kappa/16}
\] 允许换成 \(s=-4\)；原实际输入、\(A,p,a^2\) 和 \(\chi/A\) 被相应保留。

在这些代表上，标量预算给 \(p/V\le1/(2\pi^2)\)。Gauss--Bonnet 给 \[
 \fint|W^-|^2
 \le \frac{8\pi^2\chi}{V}-\frac23+\frac{P}{2V}
 \le4B-\frac23+o(1).
\] 也可使用端点更锐的 \(16B/9-2/3+o(1)\)，但此处只需要有限性。故得到 \eqref{eq:r3:1}，与上文对偶测试的矛盾。这证明对每个固定有限 \(B\) 的 \(c(B)>0\)。

\textbf{无界版本的准确剩余状态。} 若完全不控制 \(\chi/A\) 的 O10 仍有逃逸序列，则本文强迫 \[
 \fint|W^-|^2\longrightarrow\infty;
\] 否则存在有界子列，再次违反 \eqref{eq:r3:1}。由端点 \(p/V\to2/(9\pi^2)\)、\(W^+\to0\) 及 Euler 恒等式，它还须满足 \[
 \frac\chi p\longrightarrow\infty,
 \qquad \frac\chi A\longrightarrow\infty.
\] 本文并未证明这种无界密度逃逸不可能。有限 \(K\) 被用于 \eqref{eq:r3:5}、第一延拓的 adjoint 控制、\(c\) 的 \(L^4\) 和导数预算；不能把这项依赖隐去后宣布无条件 O10 已解。

\section{两个算子的区分}\label{ch3:ux4e24ux4e2aux7b97ux5b50ux7684ux533aux5206}

\(t(v)\in S^+\otimes L^4\) 的旋量 Lichnerowicz 势与 \(\Phi\wedge t(v)\in L^5\) 的线丛曲率收缩不是同一个量。原计算中前者在所用方向趋于 \(-7\)，后者给出上文的 \(+10/3\)。本章使用线丛弱配对及其合法测试，不从前一个负势推导消失。

O10 的极限幅度为固定正数；O11 位于允许的 KE 端，目标尺度是 \(Q=Ah=o(V)\)。本章的 \(o(V)\) 对偶结论不构成 O11 的 \(o(Q)\) 锐利估计。

\begin{thebibliography}{P99}
\small
\bibitem[R03]{R03} 本会话研究原稿，\emph{第五幂线丛与任意固定有限密度端点}，2026-09-29。\path{O10_Finite_Density_Fifth_Line.md}。

\bibitem[P03]{Strong} 所提供网页续研，\emph{O10 Strong Current}，2026-09-29，R1--R2 及同解曲率拟合。

\bibitem[P04]{Prolongation} 所提供网页续研 B，\emph{O10 Resolvent Prolongation}，2026-09-29，第一延拓与唯一剩余分量。

\bibitem[P05]{FD} 所提供网页续研，\emph{O10 closes at every fixed finite Euler density}，2026-09-29，\path{O10_All_Finite_Density.md}。

\end{thebibliography}
\end{document}
